\IfFileExists{stacks-project.cls}{%
\documentclass{stacks-project}
}{%
\documentclass{amsart}
}

% For dealing with references we use the comment environment
\usepackage{verbatim}
\newenvironment{reference}{\comment}{\endcomment}
%\newenvironment{reference}{}{}
\newenvironment{slogan}{\comment}{\endcomment}
\newenvironment{history}{\comment}{\endcomment}

% For commutative diagrams we use Xy-pic
\usepackage[all]{xy}

% We use 2cell for 2-commutative diagrams.
\xyoption{2cell}
\UseAllTwocells

% We use multicol for the list of chapters between chapters
\usepackage{multicol}

% This is generally recommended for better output
\usepackage{lmodern}
\usepackage[T1]{fontenc}

% For cross-file-references
\usepackage{xr-hyper}

% Package for hypertext links:
\usepackage{hyperref}

% For any local file, say "hello.tex" you want to link to please
% use \externaldocument[hello-]{hello}
\externaldocument[introduction-]{introduction}
\externaldocument[conventions-]{conventions}
\externaldocument[sets-]{sets}
\externaldocument[categories-]{categories}
\externaldocument[topology-]{topology}
\externaldocument[sheaves-]{sheaves}
\externaldocument[sites-]{sites}
\externaldocument[stacks-]{stacks}
\externaldocument[fields-]{fields}
\externaldocument[algebra-]{algebra}
\externaldocument[brauer-]{brauer}
\externaldocument[homology-]{homology}
\externaldocument[derived-]{derived}
\externaldocument[simplicial-]{simplicial}
\externaldocument[more-algebra-]{more-algebra}
\externaldocument[smoothing-]{smoothing}
\externaldocument[modules-]{modules}
\externaldocument[sites-modules-]{sites-modules}
\externaldocument[injectives-]{injectives}
\externaldocument[cohomology-]{cohomology}
\externaldocument[sites-cohomology-]{sites-cohomology}
\externaldocument[dga-]{dga}
\externaldocument[dpa-]{dpa}
\externaldocument[sdga-]{sdga}
\externaldocument[hypercovering-]{hypercovering}
\externaldocument[schemes-]{schemes}
\externaldocument[constructions-]{constructions}
\externaldocument[properties-]{properties}
\externaldocument[morphisms-]{morphisms}
\externaldocument[coherent-]{coherent}
\externaldocument[divisors-]{divisors}
\externaldocument[limits-]{limits}
\externaldocument[varieties-]{varieties}
\externaldocument[topologies-]{topologies}
\externaldocument[descent-]{descent}
\externaldocument[perfect-]{perfect}
\externaldocument[more-morphisms-]{more-morphisms}
\externaldocument[flat-]{flat}
\externaldocument[groupoids-]{groupoids}
\externaldocument[more-groupoids-]{more-groupoids}
\externaldocument[etale-]{etale}
\externaldocument[chow-]{chow}
\externaldocument[intersection-]{intersection}
\externaldocument[pic-]{pic}
\externaldocument[weil-]{weil}
\externaldocument[adequate-]{adequate}
\externaldocument[dualizing-]{dualizing}
\externaldocument[duality-]{duality}
\externaldocument[discriminant-]{discriminant}
\externaldocument[derham-]{derham}
\externaldocument[local-cohomology-]{local-cohomology}
\externaldocument[algebraization-]{algebraization}
\externaldocument[curves-]{curves}
\externaldocument[resolve-]{resolve}
\externaldocument[models-]{models}
\externaldocument[functors-]{functors}
\externaldocument[equiv-]{equiv}
\externaldocument[pione-]{pione}
\externaldocument[etale-cohomology-]{etale-cohomology}
\externaldocument[proetale-]{proetale}
\externaldocument[relative-cycles-]{relative-cycles}
\externaldocument[more-etale-]{more-etale}
\externaldocument[trace-]{trace}
\externaldocument[crystalline-]{crystalline}
\externaldocument[spaces-]{spaces}
\externaldocument[spaces-properties-]{spaces-properties}
\externaldocument[spaces-morphisms-]{spaces-morphisms}
\externaldocument[decent-spaces-]{decent-spaces}
\externaldocument[spaces-cohomology-]{spaces-cohomology}
\externaldocument[spaces-limits-]{spaces-limits}
\externaldocument[spaces-divisors-]{spaces-divisors}
\externaldocument[spaces-over-fields-]{spaces-over-fields}
\externaldocument[spaces-topologies-]{spaces-topologies}
\externaldocument[spaces-descent-]{spaces-descent}
\externaldocument[spaces-perfect-]{spaces-perfect}
\externaldocument[spaces-more-morphisms-]{spaces-more-morphisms}
\externaldocument[spaces-flat-]{spaces-flat}
\externaldocument[spaces-groupoids-]{spaces-groupoids}
\externaldocument[spaces-more-groupoids-]{spaces-more-groupoids}
\externaldocument[bootstrap-]{bootstrap}
\externaldocument[spaces-pushouts-]{spaces-pushouts}
\externaldocument[spaces-chow-]{spaces-chow}
\externaldocument[groupoids-quotients-]{groupoids-quotients}
\externaldocument[spaces-more-cohomology-]{spaces-more-cohomology}
\externaldocument[spaces-simplicial-]{spaces-simplicial}
\externaldocument[spaces-duality-]{spaces-duality}
\externaldocument[formal-spaces-]{formal-spaces}
\externaldocument[restricted-]{restricted}
\externaldocument[spaces-resolve-]{spaces-resolve}
\externaldocument[formal-defos-]{formal-defos}
\externaldocument[defos-]{defos}
\externaldocument[cotangent-]{cotangent}
\externaldocument[examples-defos-]{examples-defos}
\externaldocument[algebraic-]{algebraic}
\externaldocument[examples-stacks-]{examples-stacks}
\externaldocument[stacks-sheaves-]{stacks-sheaves}
\externaldocument[criteria-]{criteria}
\externaldocument[artin-]{artin}
\externaldocument[quot-]{quot}
\externaldocument[stacks-properties-]{stacks-properties}
\externaldocument[stacks-morphisms-]{stacks-morphisms}
\externaldocument[stacks-limits-]{stacks-limits}
\externaldocument[stacks-cohomology-]{stacks-cohomology}
\externaldocument[stacks-perfect-]{stacks-perfect}
\externaldocument[stacks-introduction-]{stacks-introduction}
\externaldocument[stacks-more-morphisms-]{stacks-more-morphisms}
\externaldocument[stacks-geometry-]{stacks-geometry}
\externaldocument[moduli-]{moduli}
\externaldocument[moduli-curves-]{moduli-curves}
\externaldocument[examples-]{examples}
\externaldocument[exercises-]{exercises}
\externaldocument[guide-]{guide}
\externaldocument[desirables-]{desirables}
\externaldocument[coding-]{coding}
\externaldocument[obsolete-]{obsolete}
\externaldocument[fdl-]{fdl}
\externaldocument[index-]{index}

% Theorem environments.
%
\theoremstyle{plain}
\newtheorem{theorem}[subsection]{Theorem}
\newtheorem{proposition}[subsection]{Proposition}
\newtheorem{lemma}[subsection]{Lemma}

\theoremstyle{definition}
\newtheorem{definition}[subsection]{Definition}
\newtheorem{example}[subsection]{Example}
\newtheorem{exercise}[subsection]{Exercise}
\newtheorem{situation}[subsection]{Situation}

\theoremstyle{remark}
\newtheorem{remark}[subsection]{Remark}
\newtheorem{remarks}[subsection]{Remarks}

\numberwithin{equation}{subsection}

% Macros
%
\def\lim{\mathop{\mathrm{lim}}\nolimits}
\def\colim{\mathop{\mathrm{colim}}\nolimits}
\def\Spec{\mathop{\mathrm{Spec}}}
\def\Hom{\mathop{\mathrm{Hom}}\nolimits}
\def\Ext{\mathop{\mathrm{Ext}}\nolimits}
\def\SheafHom{\mathop{\mathcal{H}\!\mathit{om}}\nolimits}
\def\SheafExt{\mathop{\mathcal{E}\!\mathit{xt}}\nolimits}
\def\Sch{\mathit{Sch}}
\def\Mor{\mathop{\mathrm{Mor}}\nolimits}
\def\Ob{\mathop{\mathrm{Ob}}\nolimits}
\def\Sh{\mathop{\mathit{Sh}}\nolimits}
\def\NL{\mathop{N\!L}\nolimits}
\def\CH{\mathop{\mathrm{CH}}\nolimits}
\def\proetale{{pro\text{-}\acute{e}tale}}
\def\etale{{\acute{e}tale}}
\def\QCoh{\mathit{QCoh}}
\def\Ker{\mathop{\mathrm{Ker}}}
\def\Im{\mathop{\mathrm{Im}}}
\def\Coker{\mathop{\mathrm{Coker}}}
\def\Coim{\mathop{\mathrm{Coim}}}

% Boxtimes
%
\DeclareMathSymbol{\boxtimes}{\mathbin}{AMSa}{"02}

%
% Macros for moduli stacks/spaces
%
\def\QCohstack{\mathcal{QC}\!\mathit{oh}}
\def\Cohstack{\mathcal{C}\!\mathit{oh}}
\def\Spacesstack{\mathcal{S}\!\mathit{paces}}
\def\Quotfunctor{\mathrm{Quot}}
\def\Hilbfunctor{\mathrm{Hilb}}
\def\Curvesstack{\mathcal{C}\!\mathit{urves}}
\def\Polarizedstack{\mathcal{P}\!\mathit{olarized}}
\def\Complexesstack{\mathcal{C}\!\mathit{omplexes}}
% \Pic is the operator that assigns to X its picard group, usage \Pic(X)
% \Picardstack_{X/B} denotes the Picard stack of X over B
% \Picardfunctor_{X/B} denotes the Picard functor of X over B
\def\Pic{\mathop{\mathrm{Pic}}\nolimits}
\def\Picardstack{\mathcal{P}\!\mathit{ic}}
\def\Picardfunctor{\mathrm{Pic}}
\def\Deformationcategory{\mathcal{D}\!\mathit{ef}}
\usepackage{cleveref}
\crefname{originalreference}{original reference}{original references}
\Crefname{originalreference}{Original reference}{Original references}
\makeatletter
\newlabel{categories-lemma-finite-limits-exist}{{002O}{source}{}{sourcepointer-1}{}}
\newlabel{categories-lemma-finite-limits-exist@cref}{{[originalreference][0][]002O}{[1][0][]source}{}{}{}}
\makeatother
\crefname{originallemma}{original lemma}{original lemmas}
\Crefname{originallemma}{Original lemma}{Original lemmas}
\makeatletter
\newlabel{lemma-fibre-product-affine-schemes}{{01I4}{source}{}{sourcepointer-2}{}}
\newlabel{lemma-fibre-product-affine-schemes@cref}{{[originallemma][0][]01I4}{[1][0][]source}{}{}{}}
\makeatother
\crefname{originallemma}{original lemma}{original lemmas}
\Crefname{originallemma}{Original lemma}{Original lemmas}
\makeatletter
\newlabel{lemma-glue}{{01JB}{source}{}{sourcepointer-3}{}}
\newlabel{lemma-glue@cref}{{[originallemma][0][]01JB}{[1][0][]source}{}{}{}}
\makeatother
\crefname{originallemma}{original lemma}{original lemmas}
\Crefname{originallemma}{Original lemma}{Original lemmas}
\makeatletter
\newlabel{lemma-glue-functors}{{01JJ}{source}{}{sourcepointer-4}{}}
\newlabel{lemma-glue-functors@cref}{{[originallemma][0][]01JJ}{[1][0][]source}{}{}{}}
\makeatother
\crefname{originallemma}{original lemma}{original lemmas}
\Crefname{originallemma}{Original lemma}{Original lemmas}
\makeatletter
\newlabel{lemma-morphism-into-affine}{{01I1}{source}{}{sourcepointer-5}{}}
\newlabel{lemma-morphism-into-affine@cref}{{[originallemma][0][]01I1}{[1][0][]source}{}{}{}}
\makeatother
\title{The category of schemes has finite limits}
\author{The Stacks Project Authors}
\date{stacks-project:schemes.tex:local-snapshot-20261012; GNU FDL 1.2 or later}
\begin{document}
\maketitle
\noindent\textit{Editorial note.} Counted span is the author's own proof environment, source lines 3233-3335, following the statement at 3226-3231.\par\medskip
\noindent\textit{Editorial note.} The proof opens with 'Please skip this proof', a pedagogical aside to readers, not a completeness disclaimer; the verification of (2)(a)-(2)(c) is complete.\par\medskip
\noindent\textit{Editorial note.} Condition (2)(a) is concluded by the named affine fibre-product engine with the locally-ringed-spaces parenthetical reproduced.\par\medskip
\noindent\textit{Allowed standard prerequisites (editorial).}\begin{itemize}
\item Spec(Z) is final among locally ringed spaces. lemma-morphism-into-affine (same chapter, cited at \textasciitilde{}3239-3241). Reduces existence of fibre products to representability of the fibre-product functor.
\item Functor glueing framework. lemma-glue-functors (same chapter, card L26382 this batch, cited at \textasciitilde{}3252-3253). The engine whose sheaf property and conditions (2)(a)-(2)(c) are verified for F.
\item Mapping property of the glued locally ringed space. lemma-glue (same chapter, statement at 2761-2798, cited at \textasciitilde{}3271). Justifies glueing the morphisms a\_i into a: T -> X.
\item Fibre products of affine schemes. lemma-fibre-product-affine-schemes (same chapter, statement at 1202, cited at \textasciitilde{}3310). Engine for condition (2)(a); the parenthetical notes affine fibre products are fibre products in all locally ringed spaces.
\end{itemize}\medskip
\ifdefined\hypertarget\hypertarget{sourcepointer-1}{}\fi
\noindent\textit{Source reference.} Original reference 002O: focused-seeds/geometry/S9-schemes-fibre-products/ref-tags.txt, source lines 5--5. Navigation pointer only: the counted proof cites this statement by \textbackslash{}ref as a standing input; the deduction the counted proof has to make is carried out inside the counted span, and no part of that deduction is delegated to this pointer.\par\medskip
\ifdefined\hypertarget\hypertarget{sourcepointer-2}{}\fi
\noindent\textit{Source reference.} Original lemma 01I4: sources/stacks-schemes/pinned-original-source.tex, source lines 1202--1202. Navigation pointer only: the counted proof cites this statement by \textbackslash{}ref as a standing input; the deduction the counted proof has to make is carried out inside the counted span, and no part of that deduction is delegated to this pointer.\par\medskip
\ifdefined\hypertarget\hypertarget{sourcepointer-3}{}\fi
\noindent\textit{Source reference.} Original lemma 01JB: sources/stacks-schemes/pinned-original-source.tex, source lines 2762--2762. Navigation pointer only: the counted proof cites this statement by \textbackslash{}ref as a standing input; the deduction the counted proof has to make is carried out inside the counted span, and no part of that deduction is delegated to this pointer.\par\medskip
\ifdefined\hypertarget\hypertarget{sourcepointer-4}{}\fi
\noindent\textit{Source reference.} Original lemma 01JJ: sources/stacks-schemes/pinned-original-source.tex, source lines 3110--3110. Navigation pointer only: the counted proof cites this statement by \textbackslash{}ref as a standing input; the deduction the counted proof has to make is carried out inside the counted span, and no part of that deduction is delegated to this pointer.\par\medskip
\ifdefined\hypertarget\hypertarget{sourcepointer-5}{}\fi
\noindent\textit{Source reference.} Original lemma 01I1: sources/stacks-schemes/pinned-original-source.tex, source lines 1058--1058. Navigation pointer only: the counted proof cites this statement by \textbackslash{}ref as a standing input; the deduction the counted proof has to make is carried out inside the counted span, and no part of that deduction is delegated to this pointer.\par\medskip
\begin{lemma}
\label{lemma-fibre-products}
The category of schemes has a final object, products and fibre products.
In other words, the category of schemes has finite limits, see
Categories, Lemma \ref{categories-lemma-finite-limits-exist}.
\end{lemma}

\begin{proof}
Please skip this proof. It is more important to learn
how to work with the fibre product which is explained in
the next section.

\medskip\noindent
By Lemma \ref{lemma-morphism-into-affine}
the scheme $\Spec(\mathbf{Z})$ is a final
object in the category of locally ringed spaces. Thus it
suffices to prove that fibred products exist.

\medskip\noindent
Let $f : X \to S$ and $g : Y \to S$ be morphisms of schemes. We have to show
that the functor
\begin{eqnarray*}
F : \Sch^{opp} & \longrightarrow & \textit{Sets} \\
T & \longmapsto &
\Mor(T, X) \times_{\Mor(T, S)} \Mor(T, Y)
\end{eqnarray*}
is representable. We claim that Lemma \ref{lemma-glue-functors}
applies to the functor $F$. If we prove this then the lemma is proved.

\medskip\noindent
First we show that $F$ satisfies the sheaf property in the
Zariski topology. Namely, suppose that $T$ is a scheme,
$T = \bigcup_{i \in I} U_i$ is an open covering, and
$\xi_i \in F(U_i)$ such that
$\xi_i|_{U_i \cap U_j} =  \xi_j|_{U_i \cap U_j}$ for
all pairs $i, j$. By definition $\xi_i$ corresponds to
a pair $(a_i, b_i)$ where $a_i : U_i \to X$ and $b_i : U_i \to Y$
are morphisms of schemes such that $f \circ a_i = g \circ b_i$.
The glueing condition says that
$a_i|_{U_i \cap U_j} =  a_j|_{U_i \cap U_j}$
and
$b_i|_{U_i \cap U_j} =  b_j|_{U_i \cap U_j}$.
Thus by glueing the morphisms $a_i$ we obtain a morphism
of locally ringed spaces (i.e., a morphism of schemes)
$a : T \to X$ and similarly $b : T \to Y$ (see for example
the mapping property of Lemma \ref{lemma-glue}). Moreover,
on the members of an open covering the compositions
$f \circ a$ and $g \circ b$ agree. Therefore
$f \circ a = g \circ b$ and the pair $(a, b)$ defines
an element of $F(T)$ which restricts to the pairs
$(a_i, b_i)$ on each $U_i$. The sheaf condition is verified.

\medskip\noindent
Next, we construct the family of subfunctors.
Choose an open covering by open affines
$S = \bigcup\nolimits_{i \in I} U_i$.
For every $i \in I$ choose open coverings by open affines
$f^{-1}(U_i) = \bigcup\nolimits_{j \in J_i} V_j$ and
$g^{-1}(U_i) = \bigcup\nolimits_{k \in K_i} W_k$.
Note that $X = \bigcup_{i \in I} \bigcup_{j \in J_i} V_j$
is an open covering and similarly for $Y$.
For any $i \in I$ and each pair $(j, k) \in J_i \times K_i$
we have a commutative diagram
$$
\xymatrix{
    & W_k \ar[d] \ar[rd] &   \\
V_j \ar[rd] \ar[r] & U_i \ar[rd] & Y \ar[d] \\
    & X \ar[r]  & S
}
$$
where all the skew arrows are open immersions. For such a
triple we get a functor
\begin{eqnarray*}
F_{i, j, k} : \Sch^{opp} & \longrightarrow & \textit{Sets} \\
T & \longmapsto &
\Mor(T, V_j) \times_{\Mor(T, U_i)} \Mor(T, W_k).
\end{eqnarray*}
There is an obvious transformation of functors $F_{i, j, k} \to F$
(coming from the huge commutative diagram above)
which is injective, so we may think of $F_{i, j, k}$ as a subfunctor
of $F$.

\medskip\noindent
We check condition (2)(a) of Lemma \ref{lemma-glue-functors}.
This follows directly from Lemma \ref{lemma-fibre-product-affine-schemes}.
(Note that we use here that the fibre products in the category of
affine schemes are also fibre products in the whole category of locally
ringed spaces.)

\medskip\noindent
We check condition (2)(b) of Lemma \ref{lemma-glue-functors}.
Let $T$ be a scheme and let $\xi \in F(T)$. In other words,
$\xi = (a, b)$ where $a : T \to X$ and $b : T \to Y$ are
morphisms of schemes such that $f \circ a = g \circ b$.
Set $V_{i, j, k} = a^{-1}(V_j) \cap b^{-1}(W_k)$. For any
further morphism $h : T' \to T$ we have
$h^*\xi = (a \circ h, b \circ h)$. Hence we see that
$h^*\xi \in F_{i, j, k}(T')$ if and only if
$a(h(T')) \subset V_j$ and $b(h(T')) \subset W_k$.
In other words, if and only if $h(T') \subset V_{i, j, k}$.
This proves condition (2)(b).

\medskip\noindent
We check condition (2)(c) of Lemma \ref{lemma-glue-functors}.
Let $T$ be a scheme and let $\xi = (a, b) \in F(T)$ as above.
Set $V_{i, j, k} = a^{-1}(V_j) \cap b^{-1}(W_k)$ as above.
Condition (2)(c) just means that $T = \bigcup V_{i, j, k}$
which is evident. Thus the lemma is proved and fibre products
exist.
\end{proof}

\end{document}
